\scanpage{166}
Alors grâce à la suite spectrale de Leray, $\unsure{\zeta|\xi}$ définit un élément dans $E_\infty^{1,n-1}\subset E_2^{1,n-1}$, donc un
\[
 \eta\in\mathbb H^1(S,\alpha\longmapsto\mathscr H^{n-1}
 (\Omega^\alpha_{S/T}\overset{L}{\otimes}\mathbb R f_*\Omega^\cdot_{X/S})).
 \tag*{(15)}
\]
Supposons en particulier les $\Omega^1_{S/T}$ plats, \deleted{de sorte} alors on obtient
\[
 \eta\in {\unsure{\mathbb H|H}}^1_{\mathrm{DR}}(S,\mathbb R^{n-1}f_*(\Omega^\cdot_{X/S}))
 \qquad(\Omega^1_{S/T}\text{ plat}).
 \tag*{(16)}
\]
Nous avons vu $[\,]$ qu'un tel objet s'identifie à une classe d'isomorphie de torseurs \inserted{$P$} sous $\mathbb R^{n-1}f_*(\Omega^\cdot_{X/S})$ muni d'un hom
\[
 P\xrightarrow{d_P}\Omega^1_{S/T}\otimes\mathbb R^{n-1}f_*(\Omega^\cdot_{X/S})
 \tag*{(17)}
\]
satisfaisant
\[
 \left\{\begin{aligned}
 d_P(x+\xi)&=d_P(x)+d^0(\xi)
 &\quad&\left\{\begin{aligned}
 x&\in P(\cdot),\quad \deleted{\unsure{\zeta}}\xi\in\mathbb R^{n-1}f_*(\Omega^\cdot_{X/S})(\cdot)\\
 d^0&\text{ déduit de conn. de G.M.}
 \end{aligned}\right.\\
 \operatorname{Im}d_P&\subset Z^1(\Omega^\cdot_{S/T}\otimes\mathbb R^{n-1}f_*(\Omega^\cdot_{X/S})).
 \end{aligned}\right.
 \tag*{(18)}
\]

[On sait aussi, grâce à l'hypothèse de platitude sur $\Omega^1_{S/T}$, que si $S$ est de car. nulle, la donnée de $d_P$ équivaut à la donnée de
\[
 P\xrightarrow{\delta_P^\infty}P_{S/T}^\infty(\mathbb R^{n-1}f_*(\Omega^\cdot_{X/S}))^+
 \mathrel{\vcenter{\hbox{\tikz{\draw (.09,.14)--(0,0)--(.22,0)--(.12,-.16);}}}}
 \tag*{(19)}
\]
satisfaisant la condition
\[
 \delta_P^\infty(x+\zeta)=\delta_P^\infty(x)+d_{S/T}^+(\zeta)
 \qquad x\in P(\cdot),\ \zeta\in\mathbb R^{n-1}f_*(\Omega^\cdot_{X/S})(\cdot)
 \tag*{(20)}
\]
\margin{$d^+_{S/T}=d^\infty_{S/T}-\delta^\infty_{S/T}$\\
$\swarrow$\\
$\delta^\infty_{S/T}$ défini par la stratification de $\mathbb R^{n-1}f_*(\Omega^\cdot_{X/S})$ déduite de sa connexion de G.M.}
\deleted{\unsure{(21)}}
% archive sheet number: 165

